\subsection{Berechnungen in einer Transistorschaltungen}
% Alt: Aufgabe 17
\Aufgabe{
    \label{Aufg17}
    Gegeben ist die folgende Schaltung mit dem Ausgangskennlinienfeld eines Transistors.

    \begin{figure}[H]
        \centering
        \begin{subfigure}[b]{0.45\textwidth}
            \centering
            \input{Tikz/src/FigTransistorKennlinien}\alttext{Die Abbildung zeigt eine elektronischen Schaltung und ein Diagramm mit Bezug auf den Kollektorstrom \(I_\mathrm{C}\) und die Kollektor Ermitter Spannung \(U_{\mathrm{CE}}\).}
        \end{subfigure}
        \begin{subfigure}[b]{0.45\textwidth}
            \centering
            \includesvg[width=0.9\textwidth]{Bilder/Aufgaben/FigTransistorKennlinien}\alttext{Aus dem Diagramm lässt sich der Basisstrom ermitteln. Das Diagramm kann wie folgt beschrieben werden. Links ist ein Graph mit der Beschriftung auf der y-Achse „\(I_\mathrm{C}\) in \(\mathrm{mA}\)“ und auf der x-Achse „\(U_{\mathrm{CE}}\) in \(\mathrm{V}\)“. Die y-Achse reicht von \(0\) bis \(400 \ \mathrm{mA}\), die x-Achse von \(0\) bis \(3 \ \mathrm{V}\). Es sind sieben Kurven eingezeichnet, die alle bei etwa 0 Volt starten und schnell ansteigen, dann flacher verlaufen. Jede Kurve ist mit einem Wert für \(I_\mathrm{B}\) in Milliampere beschriftet: 5; 1; 1,5; 2; 2,5; 3 und oben rechts  „\(I_\mathrm{B} = 3,5 \ \mathrm{mA}\)“.}
        \end{subfigure}
        \label{fig:FigTransistorKennlinien}
    \end{figure}

    \begin{itemize}
        \item[a)] Wie groß ist der maximale Kollektorstrom des idealen Transistors (bei $R_\mathrm{CE} = 0\,\Omega$), wenn der Kollektorwiderstand $R_\mathrm{C} = 7,5\,\Omega$ beträgt und die Versorgungsspannung $U_\mathrm{V} = 3\,\mathrm{V}$ ist?
        \item[b)] Bestimmen Sie  $U_\mathrm{CE}$, $U_\mathrm{RC}$ und $I_\mathrm{C}$ des realen Transistors bei einem Basisstrom von \( I_\mathrm{B} = \mathrm{2\,mA} \). Tragen Sie dazu die Widerstandsgerade in das Kennlinienfeld ein.
        \item[c)] Für einen neuen Arbeitspunkt fallen nur noch $1,7\,\mathrm{V}$ über dem Kollektorwiderstand $R_\mathrm{C}$ ab, der Basisstrom beträgt weiterhin $2\,\mathrm{mA}$. Bestimmen Sie den neuen Widerstand $R_\mathrm{C}$.
    \end{itemize}
}


\Loesung{
    \begin{itemize}
        \item[a)] Berechnung des maximalen Kollektorstroms bei idealem Transistor (\( R_\mathrm{CE} = 0 \)):
              \begin{equation*}
                  I_\mathrm{C}  = \frac{U_\mathrm{V}}{R_\mathrm{C}} = \frac{3\,\mathrm{V}}{7,5\,\Omega} = 400\,\mathrm{mA}
              \end{equation*}


        \item[b)] Bestimmung der Spannungen und des Kollektorstroms laut Kennlinienfeld:
              \begin{figure}[H]
                  \centering
                  \includesvg[width=0.7\textwidth]{Bilder/Aufgaben/LsgTransistorKennlinien}
                  \alttext{Aus dem Diagramm lässt sich der Basisstrom ermitteln. Das Diagramm kann wie folgt beschrieben werden. Links ist ein Graph mit der Beschriftung auf der y-Achse „\(I_\mathrm{C}\) in \(\mathrm{mA}\)“ und auf der x-Achse „\(U_{\mathrm{CE}}\) in \(\mathrm{V}\)“. Die y-Achse reicht von \(0\) bis \(400 \ \mathrm{mA}\), die x-Achse von \(0\) bis \(3 \ \mathrm{V}\). Es sind sieben Kurven eingezeichnet, die alle bei etwa 0 Volt starten und schnell ansteigen, dann flacher verlaufen. Jede Kurve ist mit einem Wert für \(I_\mathrm{B}\) in Milliampere beschriftet: 5; 1; 1,5; 2; 2,5; 3 und oben rechts  „\(I_\mathrm{B} = 3,5 \ \mathrm{mA}\)“.
                  Zusätzlich wurden in das Diagramm eine schräg verlaufende Gerade sowie Hilfslinien eingezeichnet. Die fallende Gerade stellt die Lastgerade dar. Sie beginnt bei hohem \(I_\mathrm{C}\) bei \(U_\mathrm{CE} = 0 \ \mathrm{V}\) und verläuft linear bis zu \(I_\mathrm{C} = 0 \ \mathrm{mA}\) bei \(U_\mathrm{CE} = 3 \ \mathrm{V}\).
                  Der Schnittpunkt dieser Lastgeraden mit einer der Kennlinien markiert den Arbeitspunkt. Von diesem Schnittpunkt aus sind gestrichelte Hilfslinien eingezeichnet:
                  eine senkrechte Linie zur x-Achse, eine waagerechte Linie zur y-Achse. Damit lassen sich die zugehörigen Werte für \(U_\mathrm{CE}\) und \(I_\mathrm{C}\) grafisch ablesen.
                  Unterhalb der x-Achse sind zusätzlich Spannungsabschnitte markiert. Ein Abschnitt ist mit \(U_\mathrm{CE} = 1,2 \ \mathrm{V}\) beschriftet. Der verbleibende Abschnitt ist mit \(U_\mathrm{RC} = 1,8 \ \mathrm{V}\) gekennzeichnet.}
                \caption{Kennlinienfeld mit eingezeichneter Widerstandsgerade}
              \end{figure}
              \begin{gather*}
                  U_\mathrm{CE} = 1,2\,\mathrm{V} \\
                  U_\mathrm{RC} = 1,8\,\mathrm{V} \\
                  I_\mathrm{C} = \frac{U_\mathrm{RC}}{R_\mathrm{C}} = \frac{1,8\,\mathrm{V}}{7,5\,\Omega} = 240\,\mathrm{mA}
              \end{gather*}

        \item[c)] Berechnung des neuen Widerstandswerts bei gleichem $I_\mathrm{C}$ und neuer Spannung $U_\mathrm{RC} = 1,7\,\mathrm{V}$:
              \begin{equation*}
                  R_\mathrm{C} = \frac{U_\mathrm{RC}}{I_\mathrm{C}} = \frac{1,7\,\mathrm{V}}{240\,\mathrm{mA}} = 7,08\,\Omega
              \end{equation*}
    \end{itemize}

    Siehe: Abschnitt \ref{sec:ElektrischesVerhaltenBipolartransistor}
}